\documentclass[10pt,a4paper]{amsart}
\usepackage{amsmath,amssymb,amsthm,graphicx}
% Retained preamble for the Stacks Project chapter snapshot in this directory.
% The chapter's own preamble.tex is a copy of the shared Stacks Project preamble; the block
% below is sliced MECHANICALLY, line for line, from it. Nothing was added, removed or
% redefined; the wrapper class, the xr-hyper/hyperref wiring and the externaldocument list
% are omitted because the assembler supplies the document class.
%
% lines taken from preamble.tex: 3, 5, 9, 11, 12, 156, 157, 158, ... (51 lines total)
%
% package coverage restored from the chapter's own preamble (2026-10-13): xy, xy-2cell

\usepackage[all]{xy}
\xyoption{2cell}
\UseAllTwocells

\newenvironment{reference}{\comment}{\endcomment}
\newenvironment{slogan}{\comment}{\endcomment}
\newenvironment{history}{\comment}{\endcomment}
\theoremstyle{plain}
\newtheorem{theorem}[subsection]{Theorem}
\newtheorem{proposition}[subsection]{Proposition}
\newtheorem{lemma}[subsection]{Lemma}
\theoremstyle{definition}
\newtheorem{definition}[subsection]{Definition}
\newtheorem{example}[subsection]{Example}
\newtheorem{exercise}[subsection]{Exercise}
\newtheorem{situation}[subsection]{Situation}
\theoremstyle{remark}
\newtheorem{remark}[subsection]{Remark}
\newtheorem{remarks}[subsection]{Remarks}
\numberwithin{equation}{subsection}
\def\lim{\mathop{\mathrm{lim}}\nolimits}
\def\colim{\mathop{\mathrm{colim}}\nolimits}
\def\Spec{\mathop{\mathrm{Spec}}}
\def\Hom{\mathop{\mathrm{Hom}}\nolimits}
\def\Ext{\mathop{\mathrm{Ext}}\nolimits}
\def\SheafHom{\mathop{\mathcal{H}\!\mathit{om}}\nolimits}
\def\SheafExt{\mathop{\mathcal{E}\!\mathit{xt}}\nolimits}
\def\Sch{\mathit{Sch}}
\def\Mor{\mathop{\mathrm{Mor}}\nolimits}
\def\Ob{\mathop{\mathrm{Ob}}\nolimits}
\def\Sh{\mathop{\mathit{Sh}}\nolimits}
\def\NL{\mathop{N\!L}\nolimits}
\def\CH{\mathop{\mathrm{CH}}\nolimits}
\def\proetale{{pro\text{-}\acute{e}tale}}
\def\etale{{\acute{e}tale}}
\def\QCoh{\mathit{QCoh}}
\def\Ker{\mathop{\mathrm{Ker}}}
\def\Im{\mathop{\mathrm{Im}}}
\def\Coker{\mathop{\mathrm{Coker}}}
\def\Coim{\mathop{\mathrm{Coim}}}
\DeclareMathSymbol{\boxtimes}{\mathbin}{AMSa}{"02}
\def\QCohstack{\mathcal{QC}\!\mathit{oh}}
\def\Cohstack{\mathcal{C}\!\mathit{oh}}
\def\Spacesstack{\mathcal{S}\!\mathit{paces}}
\def\Quotfunctor{\mathrm{Quot}}
\def\Hilbfunctor{\mathrm{Hilb}}
\def\Curvesstack{\mathcal{C}\!\mathit{urves}}
\def\Polarizedstack{\mathcal{P}\!\mathit{olarized}}
\def\Complexesstack{\mathcal{C}\!\mathit{omplexes}}
\def\Pic{\mathop{\mathrm{Pic}}\nolimits}
\def\Picardstack{\mathcal{P}\!\mathit{ic}}
\def\Picardfunctor{\mathrm{Pic}}
\def\Deformationcategory{\mathcal{D}\!\mathit{ef}}
\usepackage{cleveref}
\crefname{originallemma}{original lemma}{original lemmas}
\Crefname{originallemma}{Original lemma}{Original lemmas}
\makeatletter
\newlabel{lemma-lift-complex}{{0DIT}{source}{}{sourcepointer-1}{}}
\newlabel{lemma-lift-complex@cref}{{[originallemma][0][]0DIT}{[1][0][]source}{}{}{}}
\makeatother
\title{Stacks, Deformation Theory: lemma-canonical-class}
\author{The Stacks Project Authors}
\date{stacks-project:defos.tex:local-snapshot-20261012; GNU FDL 1.2 or later}
\begin{document}
\maketitle
\noindent\textit{Editorial note.} Statement span is lines 5407-5429 of sources/stacks-defos/pinned-original-source.tex; proof span is lines 5431-5527, the complete author proof (95 body lines).\par\medskip
\noindent\textit{Editorial note.} Navigation stubs: 1 cited-but-not-included labels are supplied as non-load-bearing pointers with their pinned tags where the snapshot tag index has them.\par\medskip
\noindent\textit{Editorial note.} extra\_packages is empty on purpose: the retained preamble already loads xy with options and a second declaration clashes in the assembler wrapper.\par\medskip
\noindent\textit{Editorial note.} Text is retained verbatim, including the author's own typos and placeholders.\par\medskip
\noindent\textit{Editorial note.} Queue index 6 of queue\_defos.json (pb=3630, stmt 5407-5429, proof 5431-5527).\par\medskip
\noindent\textit{Allowed standard prerequisites (editorial).}\begin{itemize}
\item Explicit prerequisite boundary. The cited inputs are prerequisites; the argument in the retained span is not replaced by any cited result.
\end{itemize}\medskip
\ifdefined\hypertarget\hypertarget{sourcepointer-1}{}\fi
\noindent\textit{Source reference.} Original lemma 0DIT: sources/stacks-defos/pinned-original-source.tex, source lines 5268--5276. Cited by the retained statement or proof but not part of the collected spans; supplied only as a navigation pointer. The argument does not depend on its content.\par\medskip
\begin{lemma}
\label{lemma-canonical-class}
Let $\mathcal{C}$ be a site. Let $\mathcal{O} \to \mathcal{O}_0$
be a surjection of sheaves of rings whose kernel is an ideal sheaf
$\mathcal{I}$ of square zero. For every object
$K_0$ in $D^-(\mathcal{O}_0)$ there is a canonical map
$$
\omega(K_0) :
K_0 \longrightarrow
K_0 \otimes_{\mathcal{O}_0}^\mathbf{L} \mathcal{I}[2]
$$
in $D(\mathcal{O}_0)$ such that for any map
$K_0 \to L_0$ in $D^-(\mathcal{O}_0)$ the diagram
$$
\xymatrix{
K_0 \ar[d] \ar[rr]_-{\omega(K_0)} & &
(K_0 \otimes^\mathbf{L}_{\mathcal{O}_0} \mathcal{I})[2] \ar[d] \\
L_0 \ar[rr]^-{\omega(L_0)} & &
(L_0 \otimes^\mathbf{L}_{\mathcal{O}_0} \mathcal{I})[2]
}
$$
commutes.
\end{lemma}

\begin{proof}
Represent $K_0$ by any complex
$\mathcal{K}_0^\bullet$ of $\mathcal{O}_0$-modules.
Apply Lemma \ref{lemma-lift-complex}
with $\mathcal{G}^n = 0$ for all $n$.
Denote $d : \mathcal{F}^n \to \mathcal{F}^{n + 1}$
the maps produced by the lemma. Then we see that
$d \circ d : \mathcal{F}^n \to \mathcal{F}^{n + 2}$
is zero modulo $\mathcal{I}$. Since $\mathcal{F}^n$ is flat,
we see that
$\mathcal{I}\mathcal{F}^n =
\mathcal{F}^n \otimes_{\mathcal{O}} \mathcal{I} =
\mathcal{F}^n_0 \otimes_{\mathcal{O}_0} \mathcal{I}$.
Hence we obtain a canonical map of complexes
$$
d \circ d : \mathcal{F}_0^\bullet
\longrightarrow
(\mathcal{F}_0^\bullet \otimes_{\mathcal{O}_0} \mathcal{I})[2]
$$
Since $\mathcal{F}_0^\bullet$ is a bounded above complex
of flat $\mathcal{O}_0$-modules, it is K-flat and may be used
to compute derived tensor product. Moreover, the map of
complexes $\mathcal{F}_0^\bullet \to \mathcal{K}_0^\bullet$
is a quasi-isomorphism by construction. Therefore the source and
target of the map just constructed represent $K_0$ and
$K_0 \otimes_{\mathcal{O}_0}^\mathbf{L} \mathcal{I}[2]$
and we obtain our map $\omega(K_0)$.

\medskip\noindent
Let us show that this procedure is compatible with maps of complexes.
Namely, let $\mathcal{L}_0^\bullet$ represent another object of
$D^-(\mathcal{O}_0)$ and suppose that
$$
\mathcal{K}_0^\bullet \longrightarrow \mathcal{L}_0^\bullet
$$
is a map of complexes. Apply Lemma \ref{lemma-lift-complex}
for the complex $\mathcal{L}_0^\bullet$, the flat modules
$\mathcal{F}^n$, the maps $\mathcal{F}^n \to \mathcal{F}^{n + 1}$,
and the compositions
$\mathcal{F}^n \to \mathcal{K}_0^n \to \mathcal{L}_0^n$
(we apologize for the reversal of letters used).
We obtain flat modules $\mathcal{G}^n$, maps
$\mathcal{F}^n \to \mathcal{G}^n$, maps
$\mathcal{G}^n \to \mathcal{G}^{n + 1}$, and maps
$\mathcal{G}^n \to \mathcal{L}_0^n$ with all properties
as in the lemma. Then it is clear that
$$
\xymatrix{
\mathcal{F}_0^\bullet \ar[d] \ar[r] &
(\mathcal{F}_0^\bullet \otimes_{\mathcal{O}_0} \mathcal{I})[2] \ar[d] \\
\mathcal{G}_0^\bullet \ar[r] &
(\mathcal{G}_0^\bullet \otimes_{\mathcal{O}_0} \mathcal{I})[2]
}
$$
is a commutative diagram of complexes.

\medskip\noindent
To see that $\omega(K_0)$ is well defined, suppose that we have two complexes
$\mathcal{K}_0^\bullet$ and $(\mathcal{K}'_0)^\bullet$
of $\mathcal{O}_0$-modules representing $K_0$ and two systems
$(\mathcal{F}^n, d : \mathcal{F}^n \to \mathcal{F}^{n + 1},
\mathcal{F}^n \to \mathcal{K}_0^n)$
and
$((\mathcal{F}')^n, d : (\mathcal{F}')^n \to (\mathcal{F}')^{n + 1},
(\mathcal{F}')^n \to \mathcal{K}_0^n)$
as above. Then we can choose a complex $(\mathcal{K}''_0)^\bullet$
and quasi-isomorphisms
$\mathcal{K}_0^\bullet  \to (\mathcal{K}''_0)^\bullet$
and
$(\mathcal{K}'_0)^\bullet  \to (\mathcal{K}''_0)^\bullet$
realizing the fact that both complexes represent $K_0$ in the
derived category. Next, we apply the result of the previous paragraph
to
$$
(\mathcal{K}_0)^\bullet \oplus (\mathcal{K}'_0)^\bullet
\longrightarrow
(\mathcal{K}''_0)^\bullet
$$
This produces a commutative diagram
$$
\xymatrix{
\mathcal{F}_0^\bullet \oplus (\mathcal{F}'_0)^\bullet
\ar[d] \ar[r] &
(\mathcal{F}_0^\bullet \otimes_{\mathcal{O}_0} \mathcal{I})[2] \oplus
((\mathcal{F}'_0)^\bullet \otimes_{\mathcal{O}_0} \mathcal{I})[2] \ar[d] \\
\mathcal{G}_0^\bullet \ar[r] &
(\mathcal{G}_0^\bullet \otimes_{\mathcal{O}_0} \mathcal{I})[2]
}
$$
Since the vertical arrows give quasi-isomorphisms on the summands
we conclude the desired commutativity in $D(\mathcal{O}_0)$.

\medskip\noindent
Having established well-definedness, the statement on compatibility
with maps is a consequence of the result in the second
paragraph.
\end{proof}

\end{document}
